%! Trigonometry and Polar Coordinates — Unit 3, Lesson 3.13 — Guided Notes (Answer Key)
\documentclass[10pt]{article}
\usepackage{precalculus-article}
\usepackage{precalculus-key}   % pulls in -boxes; do NOT also load -boxes
\newcommand{\TallMath}[1]{$\displaystyle #1\rule[-1.4em]{0pt}{3.2em}$}

\begin{document}

\pageheader{Unit 3, Lesson 3.13}{Guided Notes --- Answer Key}
\namedateperiod

\vspace{0.10in}

%----------------------------------------------------------------------
% Objectives
%----------------------------------------------------------------------
\begin{objectivebox}
\textbf{I will be able to\ldots}
\begin{enumerate}[leftmargin=1.5em, itemsep=4pt]
  \item \textbf{LO 3.13.A} --- Locate a point in the plane using polar
    coordinates $(r,\theta)$ and explain what $r$ and $\theta$ represent.
  \item \textbf{LO 3.13.A} --- Convert between polar and rectangular
    coordinates using $x = r\cos\theta$, $y = r\sin\theta$, $r = \sqrt{x^2+y^2}$,
    and $\theta = \arctan(y/x)$ with a quadrant check.
  \item \textbf{LO 3.13.A} --- Express a complex number in rectangular form
    $a+bi$ and in polar form $(r\cos\theta)+i(r\sin\theta)$.
\end{enumerate}
\end{objectivebox}

\vspace{0.08in}

%----------------------------------------------------------------------
% Vocabulary
%----------------------------------------------------------------------
\begin{vocabbox}
\begin{description}[leftmargin=0pt, itemsep=6pt]
  \item[\termblanklong{polar coordinate system}]
    A coordinate system based on \ans{concentric circles} centered at the origin
    and \ans{rays} from the origin. Points are located by $(r,\theta)$.

  \item[\termblanklong{pole}]
    The fixed reference point, which is the \ans{origin}.

  \item[\termblanklong{polar axis}]
    The reference ray from the pole pointing in the direction of the positive
    \ans{$x$}-axis.

  \item[\termblanklong{polar coordinates $(r,\theta)$}]
    An ordered pair where $|r|$ = \ans{distance} from the pole and
    $\theta$ = measure of the angle in \ans{standard} position whose
    terminal ray passes through the point.

  \item[\termblanklong{rectangular coordinates $(x,y)$}]
    An ordered pair giving the horizontal and vertical displacement from the
    origin. These are \ans{unique} for each point.
\end{description}
\end{vocabbox}

\vspace{0.08in}

%----------------------------------------------------------------------
% Hook
%----------------------------------------------------------------------
\begin{hookbox}
\textbf{Entry Question:} The point $(1,1)$ is shown in the rectangular coordinate
system below.

\vspace{4pt}
\begin{center}
\begin{tikzpicture}[x=1.2cm, y=1.2cm]
  \draw[linegray, very thin] (-0.5,-0.5) grid (2.5,2.5);
  \draw[->] (-0.5,0) -- (2.6,0) node[right]{\small $x$};
  \draw[->] (0,-0.5) -- (0,2.6) node[above]{\small $y$};
  \foreach \v in {1,2} {
    \draw (\v,0.05)--(\v,-0.05) node[below,font=\scriptsize]{\v};
    \draw (0.05,\v)--(-0.05,\v) node[left,font=\scriptsize]{\v};
  }
  \fill[navy] (1,1) circle (2.5pt) node[above right,font=\small]{$(1,1)$};
  \draw[navy, dashed] (0,0)--(1,1);
  \draw[navy!60, ->] (0.5,0) arc(0:45:0.5) node[midway,right,font=\scriptsize]{$\theta$};
\end{tikzpicture}
\end{center}

\vspace{4pt}
\textbf{Q1.} Distance from the origin to $(1,1)$: \quad $r = \sqrt{1^2+1^2} = $ \ans{$\sqrt{2}$}

\vspace{4pt}
\textbf{Q2.} Angle $\theta$ that the segment makes with the positive $x$-axis:
$\theta = $ \ans{$\pi/4$}

\vspace{4pt}
\textbf{Q3.} So we could describe $(1,1)$ using the ordered pair $(r,\theta) = $\ans{$(\sqrt{2},\, \pi/4)$}.

\end{hookbox}

\vspace{0.08in}

%----------------------------------------------------------------------
% LO 3.13.A — The Polar Coordinate System
%----------------------------------------------------------------------
\begin{notesbox}{LO 3.13.A --- The Polar Coordinate System}

In the polar coordinate system, the point $(r,\theta)$ lies on the circle of
radius \ans{$|r|$} at the angle \ans{$\theta$} (measured in standard position
from the polar axis).

\vspace{6pt}
\textbf{Key fact --- Multiple Representations:} The same point can be described
in many ways.

\vspace{4pt}
\renewcommand{\arraystretch}{1.6}
\begin{tabular}{|c|l|}
\hline
\textbf{Representation} & \textbf{How it works} \\
\hline
$(r,\, \theta + 2\pi k)$ & Adding $2\pi$ to $\theta$ makes a full extra revolution --- same terminal ray. \\
\hline
$(-r,\, \theta + \pi)$ & Negating $r$ reflects through the pole; adding $\pi$ compensates. \\
\hline
\end{tabular}

\vspace{8pt}
\textbf{Example:} Two alternate polar representations of $(3,\pi/4)$:

\vspace{4pt}
\begin{tabular}{ll}
Representation 1: & \ans{$(3,\, \pi/4 + 2\pi) = (3,\, 9\pi/4)$} \\[6pt]
Representation 2: & \ans{$(-3,\, \pi/4 + \pi) = (-3,\, 5\pi/4)$} \\
\end{tabular}

\end{notesbox}

\vspace{0.08in}

%----------------------------------------------------------------------
% LO 3.13.A — Converting Polar to Rectangular
%----------------------------------------------------------------------
\begin{notesbox}{LO 3.13.A --- Converting Polar $\to$ Rectangular}

\textbf{Formulas:}
$$x = r\cos\theta \qquad y = r\sin\theta$$

These come from scaling the unit-circle point $(\cos\theta,\sin\theta)$ by $r$.

\vspace{8pt}
\textbf{Worked Example 1:} Convert $(4, \pi/3)$ to rectangular coordinates.

\vspace{4pt}
\begin{tabular}{ll}
$x = 4\cos(\pi/3) = 4 \cdot \ans{$1/2$} = $ \ans{$2$} &
$y = 4\sin(\pi/3) = 4 \cdot \ans{$\sqrt{3}/2$} = $ \ans{$2\sqrt{3}$} \\
\end{tabular}

\vspace{4pt}
Rectangular form: \ans{$(2,\, 2\sqrt{3})$}

\vspace{8pt}
\textbf{Worked Example 2:} Convert $(3, 3\pi/4)$ to rectangular coordinates.

\vspace{4pt}
\begin{tabular}{ll}
$x = 3\cos(3\pi/4) = 3 \cdot \ans{$-\sqrt{2}/2$} = $ \ans{$-3\sqrt{2}/2$} &
$y = 3\sin(3\pi/4) = 3 \cdot \ans{$\sqrt{2}/2$} = $ \ans{$3\sqrt{2}/2$} \\
\end{tabular}

\vspace{4pt}
Rectangular form: \ans{$(-3\sqrt{2}/2,\; 3\sqrt{2}/2)$}

\end{notesbox}

\vspace{0.08in}

%----------------------------------------------------------------------
% LO 3.13.A — Converting Rectangular to Polar
%----------------------------------------------------------------------
\begin{notesbox}{LO 3.13.A --- Converting Rectangular $\to$ Polar}

\textbf{Formulas:}
$$r = \sqrt{x^2 + y^2} \qquad \theta = \arctan\!\left(\frac{y}{x}\right) \text{ (for } x > 0\text{)}
\qquad \theta = \arctan\!\left(\frac{y}{x}\right) + \pi \text{ (for } x < 0\text{)}$$

\textbf{Important:} $\arctan$ returns values in $(-\pi/2,\, \pi/2)$ only.
Always \ans{plot} the point to determine the correct quadrant.

\vspace{8pt}
\textbf{Worked Example 1 (Q1):} Convert $(3,3)$ to polar form.

\vspace{4pt}
$r = \sqrt{3^2 + 3^2} = \sqrt{\ans{18}} = $ \ans{$3\sqrt{2}$}
\qquad $\theta = \arctan\!\left(\dfrac{3}{3}\right) = \arctan(\ans{1}) = $ \ans{$\pi/4$}

\vspace{4pt}
Polar form: \ans{$(3\sqrt{2},\; \pi/4)$}

\vspace{8pt}
\textbf{Worked Example 2 (Q3):} Convert $(-1, -1)$ to polar form.

\vspace{4pt}
$r = \sqrt{(-1)^2+(-1)^2} = $ \ans{$\sqrt{2}$}

\vspace{4pt}
$x < 0$, so $\theta = \arctan\!\left(\dfrac{-1}{-1}\right) + \pi
= \arctan(\ans{1}) + \pi
= \ans{$\pi/4$} + \pi = $ \ans{$5\pi/4$}

\vspace{4pt}
Polar form: \ans{$(\sqrt{2},\; 5\pi/4)$}

\end{notesbox}

\vspace{0.08in}

%----------------------------------------------------------------------
% LO 3.13.A — Complex Numbers in Polar Form
%----------------------------------------------------------------------
\begin{notesbox}{LO 3.13.A --- Complex Numbers in Polar Form}

A complex number $a+bi$ corresponds to the point $(a,b)$ in the \textbf{complex plane}.

\vspace{4pt}
\begin{tabular}{ll}
Modulus (distance): & $r = \sqrt{a^2+b^2}$ \\[4pt]
Argument (angle):   & $\theta = \arctan(b/a)$ (with quadrant check) \\[4pt]
Polar form:         & $(r\cos\theta) + i(r\sin\theta)$ \\
\end{tabular}

\vspace{8pt}
\textbf{Fill-in table:}

\vspace{4pt}
\renewcommand{\arraystretch}{1.8}
\begin{tabular}{|l|c|c|l|}
\hline
\textbf{Complex number} & \textbf{$r$} & \textbf{$\theta$} & \textbf{Polar form} \\
\hline
$1 + i\sqrt{3}$ & \ans{$2$} & \ans{$\pi/3$} & \ans{$2\cos(\pi/3)+2i\sin(\pi/3)$} \\
\hline
$\sqrt{2} - \sqrt{2}\,i$ & \ans{$2$} & \ans{$7\pi/4$} & \ans{$2\cos(7\pi/4)+2i\sin(7\pi/4)$} \\
\hline
$-2$ & \ans{$2$} & \ans{$\pi$} & \ans{$2\cos(\pi)+2i\sin(\pi)$} \\
\hline
\end{tabular}

\end{notesbox}

\vspace{0.08in}

%----------------------------------------------------------------------
% Practice
%----------------------------------------------------------------------
\begin{practicebox}

\textbf{Guided Practice 1.}\quad Convert each polar coordinate pair to
rectangular coordinates.

\vspace{4pt}
(a)\quad $(2, \pi/6)$: $x = $ \ans{$\sqrt{3}$}, $y = $ \ans{$1$} \hfill
(b)\quad $(5, \pi)$: $x = $ \ans{$-5$}, $y = $ \ans{$0$}

\vspace{0.10in}
\textbf{Guided Practice 2.}\quad Convert each rectangular point to polar form
($r > 0$, $0 \le \theta < 2\pi$).

\vspace{4pt}
(a)\quad $(0, 4)$: $r = $ \ans{$4$}, $\theta = $ \ans{$\pi/2$} \hfill
(b)\quad $(-3, 0)$: $r = $ \ans{$3$}, $\theta = $ \ans{$\pi$}

\vspace{0.10in}
\textbf{Guided Practice 3.}\quad Write the complex number $-1 + i$ in polar form.

\vspace{4pt}
$r = $ \ans{$\sqrt{2}$} \quad $\theta = $ \ans{$3\pi/4$} \quad
Polar form: \ans{$\sqrt{2}\cos(3\pi/4) + \sqrt{2}\,i\sin(3\pi/4)$}

\end{practicebox}

\end{document}
